%!TEX root = thesis.tex

\chapter{Theoretische Grundlagen}
\label{chapter-basics}

In diesem Abschnitt wird Allgemeines über LLMs und die Softwareentwicklung mit LLMs dargestellt.
\section{Large Language Models (LLMs)}

Large Language Models werden auf großen Textkorpora trainiert, um sprachliche Muster zu modellieren und darauf basierend Textsequenzen zu erstellen. Dabei basiert ihre Funktionsweise hauptsächlich auf der
möglichen Vorhersage des nächsten Tokens innerhalb eines gegebenen Kontextes \cite{seemann2023ki}. Solche Modelle werden in der Regel mit einer sehr großen Anzahl (Millionen bis Milliarden) Parametern ausgeführt. Die zugrunde liegende Architektur basiert dabei meist auf sogenannten Transformer-Blöcken. Das Training erfolgt typischerweise zunächst durch ein selbstüberwachtes Vortraining auf großen Textkorpora.
Darauf folgen weitere Feinabstimmungs- oder Ausrichtungsprozesse, um das Modell für spezifische Aufgaben oder Anwendungen zu optimieren \cite{zhang2023survey}.

\section{LLMs in der Softwareentwicklung}

LLMs unterstützen die Softwareentwicklung in verschiedenen Bereichen. Dazu zählen unter anderem Codegenerierung, Softwaredesign, Anforderungsanalyse,
Fehlerbehebung und -erkennung, Refactoring, Leistungsoptimierung, Dokumentation sowie Softwareanalyse.
Dabei eröffnen sie neue Möglichkeiten für Automatisierung und Effizienzsteigerung \cite{LLMinSEreviewandvision}.

Die Nutzung von LLMs in der Softwareentwicklung ist in
den letzten Jahren stark gewachsen. Dennoch bestehen weiterhin zahlreiche Herausforderungen. Dazu gehören insbesondere Qualität, Verlässlichkeit, Generalisierung und Bewertung der automatisch erzeugten Ergebnisse.
Diese Aspekte stellen wichtige Themen für die zukünftige Forschung dar \cite{LLMinSEreviewandvision}.

\section{Prompt-Engineering}

Prompt-Engineering bezeichnet die gezielte Gestaltung
von Eingaben (Prompts), um das Modell so zu steuern,
dass es möglichst präzise und qualitativ hochwertige
Ergebnisse erzeugt. Ziel ist es, durch geeignete Formulierungen die Wahrscheinlichkeit der gewünschten Token-Abfolge zu erhöhen. Gut entwickelte Prompts ermöglichen die Bearbeitung verschiedener Aufgaben. Dazu gehören unter anderem Frage-Antwort-Systeme, Textklassifikation, Sprach- oder Code-Übersetzung, Codegenerierung, Dokumentation sowie logisches Schlussfolgern.

Da unterschiedliche Modelle unterschiedlich auf
Prompts reagieren können, müssen diese häufig
modellspezifisch angepasst und optimiert werden.
Neben der Formulierung spielt auch die Modellkonfiguration
eine wichtige Rolle bei der Steuerung der Ausgabe.
Beispielsweise beeinflusst die Einstellung 
die Ausgabe des Ergebnisses \cite{boonstra2024prompt}.

\section{Vibe Coding als ergänzender Ansatz}

Im Rahmen dieser Arbeit wurde neben dem strukturierten
Prompt Engineering auch ein explorativer Ansatz
untersucht, der als \textbf{Vibe-Coding} bezeichnet wird. Dabei handelt es sich um ein Vorgehen, bei
dem der Code schrittweise (iterativ) durch unsystematische Interaktion mit dem Modell
entwickelt und verbessert wird \cite{ray2024vibecoding}.

In den von \cite{ray2024vibecoding} beschriebenen Experimenten wurde der generierte
Code zunächst grob erstellt und anschließend durch
wiederholtes Nachprompten verfeinert. Fehler wurden dabei schrittweise identifiziert und
korrigiert, während neue Funktionalitäten iterativ hinzugefügt wurden. Nach der Implementierung einzelner Features wurde das Modell zusätzlich dazu aufgefordert, entsprechende Tests abzuleiten und auszuführen.
Im Vergleich zum strukturierten Prompt-Engineering zeichnet sich Vibe-Coding durch eine geringere Anfangsspezifikation aus. Statt einer detaillierten Planung erfolgt die Entwicklung schrittweise durch Feedback und Anpassungen.
Dieses Vorgehen ermöglicht eine schnelle Prototypenentwicklung und fördert kreative Lösungsansätze.
Die Autoren berichten zudem, dass dieser Ansatz ohne eine systematische Strukturierung zu inkonsistenten Ergebnissen führen kann und die iterative Fehlerkorrektur einen erhöhten manuellen Aufwand verursacht\cite{ray2024vibecoding}.


\section{KI-gestützte Migration und Re-Engineering}

KI-gestützte Migration und Re-Engineering bezeichnen
den Einsatz von Methoden der Künstlichen Intelligenz
zur Analyse, Transformation und Modernisierung
bestehender Softwaresysteme. Ziel ist es, die Funktionalität eines Legacy-Systems zu erhalten, während Implementierung oder technische Plattform modernisiert werden.

In den letzten Jahren haben LLMs zunehmend an Bedeutung für die Migration und
Modernisierung von Legacy-Systemen gewonnen.
Dies liegt insbesondere an ihrer Geschwindigkeit,
ihrem Potenzial zum Wissenserhalt und der
Verbesserung der Codequalität \cite{agarwal2024legacy}.
Gleichzeitig entstehen dabei auch Herausforderungen
und Risiken. Dazu zählen insbesondere Probleme bei der semantischen
Korrektheit durch Halluzinationen sowie mögliche
Sicherheitsrisiken \cite{chen2021evaluatinglargelanguagemodels}.

\section{Anforderungsanalyse und Systemarchitektur}

Die Anforderungsanalyse stellt einen zentralen
Bestandteil des Softwareentwicklungsprozesses dar.
Sie dient der systematischen Erfassung und
Dokumentation der Anforderungen an ein
Softwaresystem. Ziel dieses Prozesses ist es, die Erwartungen der
Stakeholder zu identifizieren und in präzise,
nachvollziehbare und überprüfbare Spezifikationen zu
überführen.

Eine klare Definition der Anforderungen bildet die
Grundlage für die spätere Systementwicklung sowie für
die Bewertung eines Systems. Grundsätzlich werden Anforderungen in funktionale und
nicht-funktionale Anforderungen unterteilt.

Funktionale Anforderungen beschreiben die konkreten
Funktionen und Dienstleistungen, die ein System
bereitstellen muss. Dazu gehören beispielsweise Prozesse zur
Benutzerregistrierung, Datenverarbeitung oder
Buchungsfunktionen. Nicht-funktionale Anforderungen beschreiben hingegen
qualitative Eigenschaften des Systems. Dazu zählen beispielsweise Leistung, Sicherheit,
Skalierbarkeit, Zuverlässigkeit oder Wartbarkeit
\cite{leucker2024anforderungen}.

Neben den Anforderungen spielt die Systemarchitektur
eine entscheidende Rolle. Sie beschreibt die grundlegende Struktur eines Systems
sowie die Beziehungen und Interaktionen zwischen den
einzelnen Komponenten.

Die Spezifikation und Dokumentation erfolgt häufig in
Form strukturierter Dokumente wie Software
Requirements Specifications (SRS) oder
Architekturdiagrammen. Diese dienen als Referenz für Entwicklung, Wartung und
Migration eines Systems. Im Kontext der Migration oder Neuentwicklung ist die
Erhaltung bestehender Anforderungen und
Architekturprinzipien von besonderer Bedeutung.
Dazu müssen funktionale und nicht-funktionale
Anforderungen sowie strukturelle Eigenschaften
analysiert und mit der neuen Implementierung
verglichen werden.

\section{Stand der Forschung}

In den letzten Jahren hat die Forschung zunehmend den
Einsatz von Large Language Models in der
Softwareentwicklung untersucht. Moderne Modelle zeigen, dass sie verschiedene Aufgaben
des Software Engineerings automatisieren können. Dazu gehören unter anderem Codegenerierung,
Codeanalyse, Dokumentation und Fehlerbehebung.

Die aktuellen Studien zeigen, dass LLMs in mehreren Phasen
des Softwareentwicklungsprozesses eingesetzt werden
können. Beispiele sind die Implementierung von Funktionen, die
Generierung von Tests sowie die Unterstützung bei
Refactoring-Prozessen \cite{LLM4SE}. Eine umfassende Übersicht liefern Fan et al., die den
Einsatz von LLMs im Software Engineering systematisch
analysieren.

Die Autoren zeigen, dass moderne Sprachmodelle in
Bereichen wie Softwaredesign, Codegenerierung,
Programmanalyse und Wartung eingesetzt werden.
Gleichzeitig werden Herausforderungen identifiziert.
Dazu gehören insbesondere die Zuverlässigkeit,
Konsistenz sowie das Problem der Halluzinationen
\cite{LLM4SE2}.

Eine Studie \cite{puerto-etal-2024-code} zeigt, dass LLMs bei bestimmten logischen und konditionalen Denkaufgaben
bessere Ergebnisse erzielen können, wenn die Problemstellung
in Form von Programmcodes anstatt ausschließlich in natürlicher
Sprache präsentiert wird. Die Autoren argumentieren, dass die
Darstellung von Informationen als Code die im Modell vorhandenen
Fähigkeiten zur strukturierten Verarbeitung und zum schrittweisen
Schlussfolgern aktiviert. Dadurch können Zusammenhänge und
Bedingungen präziser erfasst werden, was zu einer verbesserten
Lösungsqualität bei komplexen Reasoning-Aufgaben führt.
Die Ergebnisse der Studie deuten darauf hin, dass die Wahl der
Eingaberepräsentation einen wesentlichen Einfluss auf die Leistungsfähigkeit
von LLMs haben kann.


Eine weitere Studie \cite{hou2024largelanguagemodelssoftware} analysiert 395 wissenschaftliche Arbeiten zu diesem Thema. Dabei wird gezeigt, dass LLMs bereits in vielen Bereichen eingesetzt werden. Gleichzeitig bestehen weiterhin Forschungslücken, insbesondere in den Bereichen Evaluation, Sicherheit und Integration in reale Prozesse.

Diese Ergebnisse zeigen, dass LLMs ein
großes Potenzial besitzen. Dennoch bleibt offen, inwieweit sie zuverlässig für
die Analyse bestehender Systeme und die Erstellung
neuer Systeme eingesetzt werden können.